\DSource{0141}{52}
\hypertarget{AB01-PDF0141-P01}{}
\DLine{AB01-PDF0141-body-L001}{}{}{lunaribus, in quibus necessarie locus verus Lunae loco vero Solis in ecliptica opponitur; qua}
\DLine{AB01-PDF0141-body-L002}{}{p. 79.}{re elongatio inter locum Lunae iuxta eius motum medium, et gradum gradui vero Solis oppo-}
\DLine{AB01-PDF0141-body-L003}{}{}{situm, erit inaequalitas simplex motus Lunae iuxta eius locum in epicyclo. --- Nos quoque}
\DLine{AB01-PDF0141-body-L004}{}{}{multas eclipses lunares observavimus, earum tempora et medium diligentissime investigavimus,}
\DLine{AB01-PDF0141-body-L005}{5}{}{et quantitatem huius simplicis inaequalitatis ut diximus reperimus ({\itshape b}).}
\hypertarget{AB01-PDF0141-P02}{}
\DLine{AB01-PDF0141-body-L006}{}{}{\Indent Maximam quantitatem alterius inaequalitatis, contingentis ob elongationem Lunae a Sole,}
\DLine{AB01-PDF0141-body-L007}{}{}{fere $2^\circ\qfrac{2}{3}$ attingere invenerunt (1), quibus si $5^\circ1'$ simplicis inaequalitatis addantur, $7^\circ40'$ cir-}
\DLine{AB01-PDF0141-body-L008}{}{}{citer proveniunt. Hoc evenit cum epicycli centrum sit in P, quod ab A quantitate [dimidii]}
\DLine{AB01-PDF0141-body-L009}{}{}{excentrici distat (2); tunc semidiametrus epicycli [a loco Solis] declivis, erit fere $8^p$, quae}
\DLine{AB01-PDF0141-body-L010}{10}{}{sunt sinus iis $7^\circ\qfrac{2}{3}$ respondens.}
\hypertarget{AB01-PDF0141-P03}{}
\DLine{AB01-PDF0141-body-L011}{}{}{\Indent Ex iis quae diximus consequitur lineam EF, ambo centra coniungentem, $10^p19'$ (3) esse;}
\DLine{AB01-PDF0141-body-L012}{}{}{idque hoc modo demonstratur. Circa centrum A, apogeum excentrici, epicyclum HG descri-}
\DLine{AB01-PDF0141-body-L013}{}{}{bamus; lineam tangentem EH, et lineam AH ducamus. Quoniam Luna in linea tangente est,}
\DLine{AB01-PDF0141-body-L014}{}{}{inaequalitas simplex tota perficitur, quam $5^\circ1'$ esse vidimus e partibus quarum quatuor an-}
\DLine{AB01-PDF0141-body-L015}{15}{}{guli recti 360 continent; eius sinus igitur $5^p15'$ (4) iuxta rationem qua semidiametrus in 60}
\DLine{AB01-PDF0141-body-L016}{}{}{partes dividitur, ut semidiametrus circuli qui in hac figura excentricum repraesentat. --- Circa}
\DLine{AB01-PDF0141-body-L017}{}{}{centrum P, perigeum excentrici, epicyclum HG describamus, et lineam tangentem EH atque}
\DLine{AB01-PDF0141-body-L018}{}{}{lineam PH ducamus. Quoniam Luna in puncto H, scilicet in linea tangente, est, ambae inae-}
\DLine{AB01-PDF0141-body-L019}{}{}{qualitates perficiuntur, quarum summa $7^\circ40'$ complectitur, et sinus summae $8^p$ fere continebit,}
\DLine{AB01-PDF0141-body-L020}{20}{}{iuxta rationem qua quatuor anguli recti $360^\circ$ efficiunt, et semidiametrus, idest EA, in 60 par-}
\DLine{AB01-PDF0141-body-L021}{}{}{tes dividitur. --- Linea PH aequalis est lineae AH, quam $5^p\qfrac{1}{4}$ esse vidimus, si $60^p$ semidia-}
\DLine{AB01-PDF0141-body-L022}{}{}{metro EA tribuantur; sed quia centrum epicycli in locum venit quo, ob proximitatem puncti E,}
\DLine{AB01-PDF0141-body-L023}{}{}{--- quod est centrum Terrae et locus veri aspectus, --- dimensio mutari debet, id quod ex}
\DLine{AB01-PDF0141-body-L024}{}{}{observatione apparet fiet $8^p$ circiter secundum rationem qua EA $60^p$ complectitur; sed iuxta}
\DLine{AB01-PDF0141-body-L025}{25}{p. 80.}{rationem qua $8^p$ in $60^p$ convertuntur, $5^p\qfrac{1}{4}$ fient $39^p22'$. Et haec est quantitas lineae EP,}
\DLine{AB01-PDF0141-body-L026}{}{}{centrum Terrae cum perigeo excentrici coniungentis. Similiter ut $8^p$ in $5^p\qfrac{1}{4}$ convertuntur,}
\DLine{AB01-PDF0141-body-L027}{}{}{$60^p$ fiunt $39^p22'$. Si linea EP, quam $39^p22'$ esse vidimus, lineae EA, $60^p$ complectenti, addatur,}
\DLine{AB01-PDF0141-body-L028}{}{}{summa $99^p22'$ erit diametrus tota excentrici, cuius dimidia pars $49^p41'$ attinget. [Qua re}
\DLine{AB01-PDF0141-body-L029}{}{}{excentricitas $10^p19'$ complectitur ex partibus quarum AE $60^p$ continet] (5).}
\hypertarget{AB01-PDF0141-P04}{}
\DLine{AB01-PDF0141-body-L030}{30}{}{\Indent Cognita semidiametro epicycli iuxta eius declivitatem (6) respectu Solis, et cognita etiam}
\DLine{AB01-PDF0141-body-L031}{}{}{excentricitate atque semidiametro excentrici, restat nobis ut arcum HG, in tertia (7) columna}
\DLine{AB01-PDF0141-body-L032}{}{}{descriptum, supputemus et exponamus quanta fiat summa aequationis simplicis et alterius in}
\DLine{AB01-PDF0141-body-L033}{}{}{locis inter duas longinquitates, eo modo qui in tabulis descriptus est in quarta et quinta (8)}
\DLine{AB01-PDF0141-body-L034}{}{}{columna; et, quod ad quartam columnam attinet, cum hi $2^\circ40'$ fient $60'$ quae in quinta co-}
\DLine{AB01-PDF0141-body-L035}{35}{}{lumna notantur, videamus quantum ex iis coadunetur, et quae sit eius ratio ad 60 (9). Haec}
\DLine{AB01-PDF0141-body-L036}{}{}{omnia, ut nunc docebimus, cognoscuntur.}
\par\vspace{6pt}\hrule height.25pt\vspace{5pt}
\noindent\begin{minipage}[t]{.48\textwidth}\vspace{0pt}\fontsize{9}{11.5}\selectfont\setlength{\GutterOffset}{0pt}
\hypertarget{AB01-PDF0141-N01}{}
\DLine{AB01-PDF0141-notes_left-L001}{}{}{\Indent (1) Ita {\itshape Almag.} V, 3 (ed. Halma t. I, p. 293--206).}
\hypertarget{AB01-PDF0141-N02}{}
\DLine{AB01-PDF0141-notes_left-L002}{}{}{\Indent (2) Plato: « punctum P, quod est egressi circuli}
\DLine{AB01-PDF0141-notes_left-L003}{}{}{longitudinis propior », i. e. perigeum excentrici. Sen-}
\DLine{AB01-PDF0141-notes_left-L004}{}{}{sus cum verbis codicis optime convenit.}
\hypertarget{AB01-PDF0141-N03}{}
\DLine{AB01-PDF0141-notes_left-L005}{}{}{\Indent (3) Plato $10^\circ18'$; demonstratio et numeri se-}
\DLine{AB01-PDF0141-notes_left-L006}{}{}{quentes reperiuntur in {\itshape Almag.} V, 4 (ed. Halma, t. I,}
\DLine{AB01-PDF0141-notes_left-L007}{}{}{p. 296--297).}
\hypertarget{AB01-PDF0141-N04}{}
\DLine{AB01-PDF0141-notes_left-L008}{}{}{\Indent (4) Minuta in cod. $25'$.}
\hypertarget{AB01-PDF0141-N05}{}
\DLine{AB01-PDF0141-notes_left-L009}{}{}{\Indent (5) Addidi, Platonem et Ptolemaeum sequens.}
\end{minipage}
\hfill\begin{minipage}[t]{.48\textwidth}\vspace{0pt}\fontsize{9}{11.5}\selectfont\setlength{\GutterOffset}{0pt}
\hypertarget{AB01-PDF0141-N06}{}
\DLine{AB01-PDF0141-notes_right-L001}{}{}{\Indent (6) I. e. iuxta elongationem centri epicycli a}
\DLine{AB01-PDF0141-notes_right-L002}{}{}{Solis loco.}
\hypertarget{AB01-PDF0141-N07}{}
\DLine{AB01-PDF0141-notes_right-L003}{}{}{\Indent (7) Si columna numerorum singulorum, vel co-}
\DLine{AB01-PDF0141-notes_right-L004}{}{}{lumna Solis non computetur.}
\hypertarget{AB01-PDF0141-N08}{}
\DLine{AB01-PDF0141-notes_right-L005}{}{}{\Indent (8) Quinta et sexta.}
\hypertarget{AB01-PDF0141-N09}{}
\DLine{AB01-PDF0141-notes_right-L006}{}{}{\Indent (9) \Name{Schiaparelli} duce, interpretandum vi-}
\DLine{AB01-PDF0141-notes_right-L007}{}{}{detur: « et quoad quartam columnam, exhibetur in}
\DLine{AB01-PDF0141-notes_right-L008}{}{}{« ea pars proportionalis ad 60 sumenda de numeris}
\DLine{AB01-PDF0141-notes_right-L009}{}{}{« qui in quinta columna notantur ».}
\end{minipage}
